\subsection{Single machine scheduling}
\label{unicast-single}

In this section we consider a simple greedy algorithm for minimizing
the maximum delay factor on a single machine when requests have
non-uniform sizes.  We analyze the simple shortest-slack-first
($\sug$) algorithm which at any time $t$ schedules the request with
the shortest slack.  Recall that the slack of a request $J_i$ is $d_i
- a_i$ and that we have assumed without loss of generality that all
requests have uniform sizes.

\begin{center}
\begin{tabular}[r]{|c|}
\hline
\textbf{Algorithm}: \sug \\

\\

\begin{minipage}{13cm}
\begin{itemize}
\item Let $Q(t)$ be the set of alive requests at $t$.
\item Let $J_i$ be the request with the minimum slack among requests
  in $Q(t)$, ties broken by arrival time.
\item Preempt the current request
  and schedule $J_i$ if it is not being processed.






\end{itemize}
\end{minipage}\\\\

\hline
\end{tabular}
\end{center}


\begin{theorem}
\label{thm:ssf-single-machine} The algorithm $\sug$ is an $(1 +
\eps)$-speed ${1/\eps}$-competitive online algorithm for minimizing the maximum delay factor in the unicast setting.
\end{theorem}
\begin{proof}
  Consider an arbitrary request sequence $\sigma$ and let $\alpha^\sug$ be
  the maximum delay factor achieved by $\sug$ on $\sigma$. If $\alpha^\sug
  = 1$ there is nothing to prove, so assume that $\alpha^\sug > 1$.  Let
  $J_i$ be the request that witnesses $\alpha^\sug$, that is $\alpha^\sug = (f_i
  - a_i)/S_i$. Note that $\sug$ does not process any request with slack
  more than $S_i$ in the interval $[a_i, f_i]$. Let $t$ be the minimum
  value less than or equal to $a_i$ such that $\sug$ was busy processing only
  requests with slack at most $S_i$ in the interval $[t, f_i]$. It
  follows that $\sug$ had no requests with slack $\le S_i$ just before
  $t$. The total work that $\sug$ processed in $[t,f_i]$ on requests
  with slack less than equal to $S_i$ is $(1+\eps)(f_i-t)$ and all
  these requests arrive in the interval $[t,f_i]$. An optimal offline
  algorithm with $1$-speed can do total work of at most $(f_i-t)$ in
  the interval $[t,f_i]$ and hence the earliest time by which it can
  finish these requests is $f_i + \eps (f_i-t) \ge f_i + \eps(f_i -
  a_i)$. Since all these requests have slack at most $S_i$ and have
  arrived before $f_i$, it follows that $\alpha^* \ge \eps(f_i -
  a_i)/S_i$ where $\alpha^*$ is the maximum delay factor of the
  optimal offline algorithm with $1$-speed machine. Therefore, we have
  that $\alpha^{\sug}/\alpha^* \le 1/\eps$.
\end{proof}

\begin{remark}
  For uniform processing time requests, the algorithm that \emph{non-preemptively} schedules requests with the shortest slack is a
  $(1+\eps)$-speed $2/\eps$-competitive online  algorithm for
  minimizing the maximum delay factor.
\end{remark}

\subsection{Multiple machine scheduling}
\label{unicast-multi}

We now consider minimizing the maximum delay factor when there are $m$
machines. In the multiple machine setting, at any time a machine can
choose to process at most one request and a request can be processed
by at most one machine at a given time. A request is allowed to
migrate and be processed by different machines at different times; as
we remarked earlier, a schedule or algorithm is non-migratory if it
does not migrate any request. To adapt $\sug$ to this setting we take
intuition from previous work on minimizing $L_k$ norms of flow time
and stretch~\cite{BansalP03,AvrahamiA03,ChekuriGKK04}. We develop an
algorithm that immediately dispatches an arriving request to a
machine, and further does not migrate an assigned request to a
different machine once it is assigned. Each machine essentially runs
the single machine $\sug$ algorithm and thus the only remaining
ingredient to describe is the dispatching rule. For this purpose the
algorithm groups requests into classes based on their slack. A request
$J_i$ is said to be in class $k$ if $S_i \in [2^k, 2^{k+1})$. The
algorithm maintains the total processing time of requests (referred to
as \emph{volume}) that have been assigned to machine $x$ in each class
$k$. Let $U^x_{=k}(t)$ denote the total processing time of requests
assigned to machine $x$ by time $t$ of class $k$. With this notation,
the algorithm $\mmug$ (for $\sug$ with immediate dispatch) can be
described.

\begin{center}
\begin{tabular}[r]{|c|}
\hline
\textbf{Algorithm}: \mmug \\

\\

\begin{minipage}{13cm}
\begin{itemize}
\item When a new request $J_i$ of class $k$ arrives at time $t$,
assign it to a machine $x$ where $U^x_{=k}(t) = \min_y U^y_{=k}(t)$.
\item Use $\sug$ on each machine separately.
\end{itemize}
\end{minipage}\\\\

\hline
\end{tabular}
\end{center}

The rest of this section is devoted to the proof of the following
theorem.
\begin{theorem}
\label{thm:multiple} $\mmug$ is a $(1+ \eps)$-speed $O({1/
\eps})$-competitive online algorithm for minimizing the maximum delay
factor on $m$ identical machines.
\end{theorem}

We need a fair amount of notation. For each time $t$, machine $x$, and
class $k$ we define several quantities. For example $U^x_{=k}(t)$ is
the total volume assigned to machine $x$ in class $k$ by time $t$. We
use the predicate ``$\le k$'' to indicate classes $1$ to $k$.  Thus $U^x_{\le
  k}(t)$ is the total volume assigned to machine $x$ in classes $1$ to
$k$. We let $R^x_{=k}(t)$ to denote the remaining processing time on
machine $x$ at time $t$ and let $P^x_{=k}(t)$ denote the total volume
that $x$ has finished on requests in class $k$ by time $t$. Note that
$P^x_{=k}(t) = U^x_{=k}(t) - R^x_{=k}(t)$. All these quantities refer
to the algorithm $\mmug$. We use $V^*_{=k}(t)$ and $V_{=k}(t)$ to
denote the remaining volume of requests in class $k$ in an optimal
offline algorithm with speed $1$ and $\mmug$ with speed $(1+\eps)$,
respectively. Observe that $V_{=k}(t) = \sum_x R^x_{=k}(t)$.  The
quantities $V^*_{\le k}(t)$ and $V_{\le k}(t)$ are defined
analogously.

The algorithm $\mmug$ balances the amount of processing time for
requests with similar slack.  Note that the assignment of requests is
not based on the current volume of unfinished requests on the
machines, rather the assignment is based on the volume of requests
that were assigned in the past to different machines.  We begin our
proof by showing that for any $k$ the total volume of requests of
class at most $k$ that is assigned is balanced across the machines.
This easily follows from the dispatching policy of the
algorithm. Several of the lemmas below and their proofs
follow from the work in~\cite{AvrahamiA03,ChekuriGKK04}.

\begin{proposition}
\label{observation} 
For any time $t$ and two machines $x$ and $y$,
$|U^{x}_{=k}(t) - U^{y}_{=k}(t)| \leq 2^{k+1}$.  This also implies
that $|U^{x}_{\leq k}(t) - U^{y}_{\leq k}(t)| \leq 2^{k+2}$.
\end{proposition}
\begin{proof}
  The first inequality holds since all of the requests of class $k$
  are of size $\leq 2^{k+1}$.  The second inequality follows easily
  from the first.
\end{proof}




\begin{lemma}
  \label{lem:processbound} Consider any two machines $x$ and $y$.  At any time $t$, the
  difference in volume of requests that already have been
  processed is bounded as $|P^{x}_{\leq k} (t) - P^{y}_{\leq k} (t)|
  \leq 2^{k+2}$.
\end{lemma}

\begin{proof}
  Suppose the lemma is false. Then there is the first time $t_0$ when
  $P^{x}_{\leq k} (t_0) - P^{y}_{\leq k} (t_0) = 2^{k+2}$ and small
  constant $\delta t > 0$ such that $P^{x}_{\leq k} (t_0 + \delta t) -
  P^{y}_{\leq k} (t_0 + \delta t) > 2^{k+2}$. Let $t' = t_0+\delta t$.
  For this to occur, $x$ processes a request of class $\le k$ during
  the interval $I = [t_0, t']$ while $y$ processes a request of class
  $> k$. Since each machine uses $\sug$, it must be that $y$ had no
  requests in classes $\le k$ during $I$ which implies that $U^y_{\le
    k}(t') = P^y_{\le k}(t')$. Therefore,
\[
U^{y}_{\leq k} (t') = P^{y}_{\leq k} (t') < P^{x}_{\leq k} (t') - 2^{k+2}
\le U^{x}_{\leq k} (t') - 2^{k+2}\,,
\]
since $P^{x}_{\leq k} (t') \le U^{x}_{\leq k} (t')$. However, this implies
that
\[
U^{y}_{\leq k} (t') < U^{x}_{\leq k} (t')  - 2^{k+2}\,,
\]
a contradiction to \expref{Proposition}{observation}.
\end{proof}


\begin{lemma}
  \label{lem:closebound} At any time $t$ the difference between the
  residual volume of requests that needs to be processed, on any two
  different machines $x$ and $y$, is bounded as $|R^{x}_{\leq k} (t) -
  R^{y}_{\leq k}(t)| \leq 2^{k+3}$.
\end{lemma}

\begin{proof}
Combining \expref{Proposition}{observation}, \expref{Lemma}{lem:processbound},
and the fact that $R^{x}_{\leq k}(t) = U^{x}_{\leq k}(t) - P^{x}_{\leq k}(t)$ for any $x$ and $k$ by definition then,
\[
|R^{x}_{\leq k} (t) -  R^{y}_{\leq k}(t)| \le
|U^{x}_{\leq k} (t) -  U^{y}_{\leq k}(t)| + |P^{x}_{\leq k} (t) -  P^{y}_{\leq k}(t)| \le 2^{k+3}\,.
\]
\end{proof}


\begin{corollary}
  \label{cor:optarrival} At any time $t$, $V^*_{\leq k} (t) \geq V_{\leq
    k} (t) - m2^{k+3}$.
\end{corollary}

%
%
%
%
%

We are now ready to upper bound the competitiveness of \mmug~ when given $(1+\eps)$-speed in a similar fashion to the single machine
case. Consider an arbitrary request sequence $\sigma$ and let $J_i$ be
the request that witnesses the maximum delay factor $\alpha^\mmug$ of $\mmug$. Let $k$ be the class of $J_i$ and let $x$ be the machine $J_i$ was processed on by $\mmug$.  We know that $\alpha^\mmug = (f_i - a_i)/S_i$ by definition of $J_i$. We use $\alpha^*$ to denote the delay factor of some fixed optimal
algorithm that uses $m$ machines of speed $1$.

Let $t$ be the last time before $a_i$ when machine $x$ processed a request
of class $> k$ under $\mmug$'s schedule. Note that $t \leq a_i$ since $x$ does not process any
request of class $> k$ in the interval $[a_i,f_i]$. At time $t$ we
know by \expref{Corollary}{cor:optarrival} that $V^*_{\leq k} (t) \geq
V_{\leq k} (t) - m2^{k+3}$. If $f_i \le a_i + 2^{k+4}$ then $\mmug$
achieves a competitive ratio of $16$ since $J_i$ is in class
$k$. Thus we will assume from now on that $f_i > a_i + 2^{k+4}$.

During the interval $I = [t, f_i)$, $\mmug$ completes a total volume of
$(1 + \eps)(f_i - t)$ work on machine $x$. Using \expref{Lemma}{lem:processbound}, any
other machine $y$ also processes a volume of $(1+\eps)(f_i-t) -
2^{k+3}$ work during $I$ of requests in classes at most $k$ . Thus the total volume processed by $\mmug$ during $I$ of requests in classes at most $ k$ is at least $m(1+\eps)(f_i-t) -
m2^{k+3}$.  During $I$, the optimal algorithm schedules at most a
$m(f_i-t)$ volume of work for requests in classes at most $k$.  Combining this with
\expref{Corollary}{cor:optarrival}, we see that
\begin{eqnarray*}
V^*_{\le k}(f_i) & \ge & V_{\le k}(t) - m2^{k+3} + m(1+\eps)(f_i-t)  - m2^{k+3} \\
  & \ge & V_{\le k}(t) + m(1+\eps)(f_i-t) - m2^{k+4} \ge   \eps m (f_i - t)\,.
\end{eqnarray*}
In the last inequality we use the fact that $f_i-t \ge f_i -
a_i \ge 2^{k+4}$. Without loss of generality assume that no
requests arrives exactly at $f_i$. Therefore $V^*_{\le k}(f_i)$ is the
total volume of requests in classes $1$ to $k$ that the optimal
algorithm has left to finish at time $f_i$ and all these requests have
arrived before $f_i$. The earliest time that the optimal algorithm can
finish all these requests is $f_i + \eps (f_i-t)$ and therefore it
follows that $\alpha^* \ge \eps (f_i - t)/2^{k+1}$. Since $\alpha^\mmug \le (f_i-a_i)/2^k$ and $t \leq a_i$, it follows that $\alpha^\mmug \le 2 \alpha^*/\eps$.

Thus $\alpha^\mmug \le \max\{16, 2\alpha^*/\eps\}$ which completes the proof
of \expref{Theorem}{thm:multiple}.
